\chapter{Bonding Modes}
\label{ch:modes}

The machine has four bonding modes. All four make two welds and a loop; they
differ in who controls the descent and how the wire is torn.

You choose a mode by choosing a configuration --- each configuration belongs to
one mode. Switch modes with $\leftarrow$ and $\rightarrow$ on the configuration
list, or press \keycap{MANUAL} to jump straight to Manual.

\section{Choosing a mode}

\begin{center}\small
\begin{tabular}{@{}L{24mm}L{24mm}L{24mm}L{44mm}@{}}
\toprule
\textbf{Mode} & \textbf{Descent} & \textbf{Tear by} & \textbf{Use it for} \\
\midrule
Semi Auto  & automatic & T axis & Normal production work \\
Manual     & you drive & T axis & Setting up, awkward geometry, one-offs \\
Table Tear & automatic & Y stage & Work where the T axis cannot reach \\
Lange Cplr & automatic & T axis & Bonds needing a drag into the second weld \\
\bottomrule
\end{tabular}
\end{center}

\section{Semi Auto}

The reference mode. You trigger each of the two bonds; everything else is
automatic.

\subsection*{First bond}

\begin{enumerate}[leftmargin=*]
\item \textbf{Press and hold} the right mouse button. The machine loads
  tracking force and moves the head to \menuitem{Search 1}.
\item \textbf{Release}. The head descends toward \menuitem{Overtravel} under
  the light constant force, stopping when it touches the pad.
\item The machine waits \menuitem{Contact Stl}, ramps to \menuitem{Force 1 G},
  waits again, then finds resonance and delivers \menuitem{Energy 1} at
  \menuitem{Power 1}.
\item Force drops back to constant and the machine waits \menuitem{Cooling}.
\end{enumerate}

\subsection*{Loop}

\begin{enumerate}[leftmargin=*,resume]
\item The clamp opens and the head rises to \menuitem{Kink Height}, forming the
  kink. It waits for the wire to release.
\item The T axis feeds \menuitem{Tail} of wire.
\item The Y stage backs off by \menuitem{Reverse}.
\item The head rises to \menuitem{Loop Height} and waits for you --- for as
  long as you need. The clamp stays open while it waits, so the wire is never
  held rigid.
\end{enumerate}

\subsection*{Second bond}

\begin{enumerate}[leftmargin=*,resume]
\item \textbf{Press and hold} the right button again. The clamp pulses closed
  and open to take up slack. The head moves to \menuitem{Search 2} and the
  stage to \menuitem{Stepback}.
\item \textbf{Release}. The head descends and touches down.
\item Settle, ramp to \menuitem{Force 2 G}, settle, then \menuitem{Energy 2} at
  \menuitem{Power 2}. Cool.
\end{enumerate}

\subsection*{Tear and tail}

\begin{enumerate}[leftmargin=*,resume]
\item The clamp closes and the T axis moves \menuitem{Tear}, breaking the wire.
\item The head returns to \menuitem{Reset Height} and waits
  \menuitem{Tail Delay}.
\item The T axis returns to its origin while a burst of \menuitem{Tail Energy}
  at \menuitem{Tail Power} assists the wire feed, forming the fresh tail.
\item The stage returns to its origin. The cycle is over.
\end{enumerate}

\section{Manual}

You drive the descents yourself with the mouse buttons. Use it when the
automatic search cannot be trusted --- unusual heights, delicate parts, or
while setting a job up.

\begin{center}\small
\begin{tabular}{@{}L{26mm}L{95mm}@{}}
\toprule
\textbf{Button} & \textbf{Action} \\
\midrule
Left held  & Head moves down \\
Right held & Head moves up \\
Neither    & Head stops where it is \\
\bottomrule
\end{tabular}
\end{center}

The rest of the cycle --- the settle, the weld, the kink, the tail, the tear
--- runs automatically exactly as in Semi Auto.

Differences from Semi Auto:

\begin{itemize}[leftmargin=*]
\item \menuitem{Search 1}, \menuitem{Search 2}, \menuitem{Stepback} and
  \menuitem{Overtravel} are hidden --- you are the search.
\item The Y stage makes a single kink stroke: it backs off by
  \menuitem{Reverse} at the very start of the cycle and returns to its origin
  at the kink, rather than backing off after the kink.
\item There is no clamp slack pulse before the second bond.
\end{itemize}

\begin{wbwarning}
In Manual mode nothing limits how far down you drive the head except the end of
its travel. Lower it gently and watch the work.
\end{wbwarning}

\section{Table Tear}

Identical to Semi Auto up to the end of the second weld. The T axis then stays
still and the Y stage forms and tears the tail instead.

\begin{enumerate}[leftmargin=*]
\item After the second weld, the head moves to \menuitem{Second Z Hgt} and
  waits for the wire to release.
\item The stage moves to \menuitem{Table Tail}, forming the tail.
\item The clamp closes.
\item The stage moves to \menuitem{Table Tear}, tearing the wire.
\item The machine waits \menuitem{Tear Stabil}, then the stage and head return
  together.
\end{enumerate}

\begin{wbnote}
\menuitem{Table Tail} and \menuitem{Table Tear} are absolute stage positions,
unlike \menuitem{Tail} and \menuitem{Tear} which are relative travels. Setting
them is therefore more direct: the numbers mean where the stage goes, not how
far it moves.
\end{wbnote}

There is no ultrasonic tail assist in this mode --- ultrasonics are used only
for the two welds.

\section{Lange Cplr}

Semi Auto with two changes, for bonds that need the tool to drag as it lands:

\begin{itemize}[leftmargin=*]
\item After the kink, the stage does \emph{not} back off. The head goes
  straight from \menuitem{Kink Height} to \menuitem{Loop Height}.
\item During the second-bond descent, the stage moves by \menuitem{Reverse}
  \emph{at the same time} as the head comes down, so the tool scrubs into
  position as it lands.
\end{itemize}

Everything else --- the triggers, the welds, the T-axis tear, the tail assist
--- is as Semi Auto.

\TODO{Describe when Lange Coupling is the right choice in this shop's work,
with an example device or geometry, and give a starting value for
\menuitem{Reverse} in that mode. The firmware defaults \menuitem{Reverse} to
zero, which makes Lange Coupling behave almost identically to Semi Auto until
somebody sets it.}

\section{What each mode shows}

The parameter screens change with the mode --- a parameter the mode cannot use
is hidden rather than shown and ignored. See Appendix~\ref{app:defaults}.
